
\subsection{Using OpenJML}

\begin{frame}{\insertsubsection}
	
	\begin{fancycolumns}
		
		\begin{definition}{Task (10 minutes)}
			
			\begin{enumerate}
				\item (Windows) Activate WSL2
				\item Download a fitting OpenJML version
				\item Unpack/prepare
				\item Test using the example \\ \texttt{./openjml -esc demos/MaxBad.java}
			\end{enumerate}
		\end{definition}
		
		\nextcolumn
		
		\begin{example}{Download}
			\centering\fancyqr{color=black,height=30mm}{https://github.com/OpenJML/OpenJML/releases/}
			
		\end{example}
		
	\end{fancycolumns}
\end{frame}

\subsection{Preconditions}

\begin{frame}{\insertsubsection}
	
	\begin{fancycolumns}
		
		\begin{definition}{Task (15 minutes)}
			
			\begin{enumerate}
				\item Skim through the OpenJML documentation on Preconditions
				\item Try to solve Question 2
				\item Test your solution with an OpenJML call
			\end{enumerate}
		\end{definition}
		
		\nextcolumn
		
		\begin{example}{Links}
			Documentation: 
			\begin{center}
				\centering\fancyqr{color=black,height=20mm}{https://www.openjml.org/tutorial/Preconditions}
			\end{center}
			
			
			Tasks: 
			\begin{center}
				\centering\fancyqr{color=black,height=20mm}{https://www.openjml.org/tutorial/exercises/PreCondEx.html}
			\end{center}
		\end{example}
		
	\end{fancycolumns}
\end{frame}

\subsection{Postconditions}

\begin{frame}{\insertsubsection}
	
	\begin{fancycolumns}
		
		\begin{definition}{Task (15 minutes)}
			
			\begin{enumerate}
				\item Skim through the OpenJML documentation on Postconditions
				\item Try to solve Question 3
				\item Test your solution with an OpenJML call
			\end{enumerate}
		\end{definition}
		
		\nextcolumn
		
		\begin{example}{Links}
			Documentation: 
			\begin{center}
				\centering\fancyqr{color=black,height=20mm}{https://www.openjml.org/tutorial/Postconditions}
			\end{center}
			
			
			Tasks: 
			\begin{center}
				\centering\fancyqr{color=black,height=20mm}{https://www.openjml.org/tutorial/exercises/PostCondEx.html}
			\end{center}
		\end{example}
		
	\end{fancycolumns}
\end{frame}

\subsection{Method Calls}

\begin{frame}{\insertsubsection}
	
	\begin{fancycolumns}
		
		\begin{definition}{Task (15 minutes)}
			
			\begin{enumerate}
				\item Skim through the OpenJML documentation on Method Calls
				\item Try to solve Question 1
				\item Test your solution with an OpenJML call
			\end{enumerate}
		\end{definition}
		
		\nextcolumn
		
		\begin{example}{Links}
			Documentation: 
			\begin{center}
				\centering\fancyqr{color=black,height=20mm}{https://www.openjml.org/tutorial/MethodCalls}
			\end{center}
			
			
			Tasks: 
			\begin{center}
				\centering\fancyqr{color=black,height=20mm}{https://www.openjml.org/tutorial/exercises/MethodCallsEx.html}
			\end{center}
		\end{example}
		
	\end{fancycolumns}
\end{frame}

\subsection{Loops}

\begin{frame}{\insertsubsection}
	
	\begin{fancycolumns}
		
		\begin{definition}{Task (20 minutes)}
			
			\begin{enumerate}
				\item Skim through the OpenJML documentation on Loops
				\item Try to solve Question 1a
				\item Test your solution with an OpenJML call
			\end{enumerate}
		\end{definition}
		
		\nextcolumn
		
		\begin{example}{Links}
			Documentation: 
			\begin{center}
				\centering\fancyqr{color=black,height=20mm}{https://www.openjml.org/tutorial/Loops}
			\end{center}
			
			
			Tasks: 
			\begin{center}
				\centering\fancyqr{color=black,height=20mm}{https://www.openjml.org/tutorial/exercises/SpecifyingLoopsEx.html}
			\end{center}
		\end{example}
		
	\end{fancycolumns}
\end{frame}
